\section{PROCEDIMENTOS METODOLÓGICOS}\label{sec:metodologia}

Três decisões de projeto decorrem do instrumento que autorizou o tratamento dos dados --- termo
firmado pela administração na condição de controladora \cite{seadpi2026termo}. A distinção é
mantida ao longo de todo o trabalho: a
\emph{obrigação}, em cada caso, é contratual e decorre daquele instrumento; a legislação fornece
os padrões que o instrumento implementa e os interesses que ele protege.

A Figura~\ref{fig:pipeline} exibe o \emph{pipeline} inteiro, e o percurso que ela desenha é o mesmo
que as subseções seguintes detalham. Da extração sobre o espelho analítico do sistema
de origem, que aplica o recorte da seção~\ref{sec:corpus}, os eventos são carregados em um grafo de
conhecimento de eventos (seção~\ref{sec:ekg}); o documento inicial de cada processo administrativo é
recuperado, convertido a texto e submetido à etapa de supressão de dados pessoais que o instrumento
autorizador exige antes de qualquer vetorização (seção~\ref{sec:documento}); dele sai também o
vocabulário controlado de palavras-chave (seção~\ref{sec:keywords}), que é a única ponte legível
entre o vetor e a prosa do catálogo. Desse material saem os seis blocos de atributos que representam
cada processo (seção~\ref{sec:atributos}), sobre os quais três famílias de agrupamento derivam os
tipos (seção~\ref{sec:agrupamento}). Do mesmo grafo sai o traço de cada processo,
renderizado nos três alfabetos da seção~\ref{sec:alfabetos}, e é sobre esses traços que se ajustam,
por tipo, os modelos de caminho típico e a decodificação restrita da seção~\ref{sec:caminhos}. Cada
tipo recebe também nome e descrição gerados a partir de seus próprios membros
(seção~\ref{sec:nomeacao}), o que fecha a entrada de catálogo. O
protocolo de avaliação (seção~\ref{sec:protocolo}), o estimando que ele produz
(seção~\ref{sec:estimando}) e a camada de evidência que tenta quebrá-lo
(seção~\ref{sec:evidencia}) fecham o percurso, e é dele que saem as tabelas e o catálogo entregues.

\textbf{Cada etapa declara a pergunta que responde, e essa pergunta é copiada para dentro de cada
artefato que a etapa escreve}, ao lado do
identificador da versão do código e do resumo criptográfico de cada entrada consumida. Um leitor
que discorde de uma conclusão pode localizar o artefato que a produziu, ler a pergunta que ele diz
responder e reexecutar a etapa sobre as mesmas entradas.

\begin{figure}[htbp]
\centering
\includegraphics[width=\textwidth]{pipeline-macro.pdf}
\caption{Etapas do \emph{pipeline} de processamento, das fontes de dados ao catálogo de tipos de
processo e de seus caminhos típicos, com os artefatos que cada etapa produz.}
\label{fig:pipeline}
\end{figure}

\subsection{Caracterização da pesquisa}\label{sec:caracterizacao}

Quanto à natureza, esta é uma \textbf{pesquisa científica aplicada}: o produto não é uma
proposição teórica sobre agrupamento de traços, e sim um artefato --- um catálogo de tipos de
processo administrativo e de seus caminhos típicos --- destinado a ser operado por uma
administração pública concreta. Quanto à abordagem, é \textbf{quantitativa}: todas as afirmações
substantivas são medidas, comparadas contra controles explícitos e acompanhadas de intervalo ou de
teste. Quanto aos objetivos, é \textbf{exploratória e descritiva} na etapa de derivação dos tipos,
uma vez que não há taxonomia prévia confiável contra a qual comparar, e \textbf{explicativa} na
etapa de avaliação, na qual se procura estabelecer qual componente do método carrega o sinal
observado. Quanto ao procedimento, é um \textbf{estudo de caso único}, na acepção corrente em
Engenharia de Software: uma organização, um sistema de origem, uma janela temporal.

O caráter de caso único decorre do instrumento autorizador. Um segundo órgão exigiria um
segundo instrumento autorizador, e o existente delimita o objeto a esta administração e ao ano de
2025 (cláusula 2.1, seção~\ref{sec:corpus}). O estudo descreve a SEAD-PI no período autorizado; a
seção 4.11 delimita o alcance.

\subsubsection*{Sujeitos}

Os \emph{sujeitos} desta pesquisa, em sentido estrito, são os processos administrativos ---
registros de tramitação, não pessoas. O cenário em que eles correm, e as três propriedades da
implantação do SEI que determinam o que o registro pode mostrar, são descritos na
seção~\ref{sec:cenario}.

\subsubsection*{Posicionalidade}

Sou servidor da SEAD-PI, atuando como analista de dados, e graduando em Engenharia de
Software. A declaração é dupla e as duas faces precisam ser lidas juntas. De um lado, foi esse
vínculo que tornou o acesso possível e que permitiu interpretar as unidades organizacionais e os
tipos de ação sem o intermédio de documentação incompleta. De outro, o mesmo vínculo é fonte
potencial de viés no enquadramento e na interpretação, e a Secretaria deteve --- por força da
cláusula 4.4.1 do instrumento autorizador --- direito de aprovação prévia sobre a divulgação
acadêmica dos resultados aqui reportados.

\subsubsection*{Convenção terminológica}

Em português, a palavra ``processo'' colide de forma particularmente danosa: designa ao mesmo tempo
o caso administrativo, o procedimento genérico e o objeto da disciplina que se aplica a ele. A
convenção abaixo é fixada aqui e mantida em todo o trabalho, sem exceção; ela também consta do
Apêndice A.

\begin{itemize}
  \item \textbf{Processo administrativo} --- o caso individual, unidade de análise deste trabalho:
    um processo gerado no SEI, com número próprio, documento inicial e histórico de tramitação. Quando o
    contexto for inequívoco, admite-se a forma abreviada ``processo''; nunca no sentido de tipo nem
    no de disciplina.
  \item \textbf{Tipo de processo} --- o agrupamento emergente, isto é, o conjunto de processos
    administrativos que o método reúne por semelhança. É o objeto que o catálogo cataloga. Não se
    confunde com o \emph{tipo de processo declarado} do SEI, que é o rótulo escolhido pelo servidor
    no momento da autuação, e cuja confiabilidade é justamente o que este trabalho põe à prova.
  \item \textbf{Mineração de processos} --- a disciplina \cite{vanderaalst2016processmining,
    ieee2012manifesto}, e apenas ela.
  \item \textbf{Traço} --- tradução de \emph{trace}: a sequência ordenada de eventos de um único
    processo administrativo, renderizada em um dos três alfabetos da seção~\ref{sec:alfabetos}.
  \item \textbf{Ação por unidade} --- o alfabeto não projetado, em que cada estado é o par
    \emph{ação@unidade}. Os outros dois alfabetos são projeções dele. Adverte-se que no código este
    alfabeto se chama \texttt{unit}, embora seja o único dos três que carrega \emph{ambos} os eixos.
  \item \textbf{Roteamento} --- tradução de \emph{routing}: o \emph{alfabeto de traço} que conserva
    o itinerário de unidades organizacionais por onde o processo administrativo passa, abstraído do
    que ocorreu dentro de cada uma.
  \item \textbf{Procedimento} --- o \emph{alfabeto de traço} que conserva a sequência de ações,
    abstraída do organograma. É a representação mais transferível dos três, por não depender dos
    nomes de unidade desta administração.
  \item \textbf{Tramitação} --- o movimento administrativo em si, isto é, o envio de um processo
    administrativo de uma unidade a outra. Os dois termos são próximos e a distinção entre eles é
    mantida sem exceção em todo o trabalho: roteamento é a representação, tramitação é o que a
    administração faz.
  \item \textbf{Registros de atividades} --- tradução de \emph{event log}: o conjunto das
    atividades que o sistema de origem registrou, isto é, a base sobre a qual todo o método
    opera. Prefere-se essa forma a ``registro de atividades de eventos'' por ser a que a própria administração
    emprega ao se referir ao que o SEI armazena.
  \item \textbf{Ação} --- tradução de \emph{action}: o tipo de andamento registrado pelo sistema
    (gerar documento, assinar, concluir, enviar). \textbf{Atividade} designa a ocorrência
    individual de uma ação em um processo administrativo, isto é, a linha dos registros de atividades.
\end{itemize}

\subsection{Cenário da pesquisa: a SEAD-PI e o processo eletrônico}\label{sec:cenario}

A Secretaria de Administração do Estado do Piauí é o órgão central dos sistemas de pessoal e de
administração geral do Poder Executivo estadual. Sua carga de trabalho é predominantemente
\emph{interna}: licenças, cessões, estágio, consignações e descontos em folha, contratações e
matérias correlatas. Parte dos procedimentos atravessa outras secretarias, outros entes federativos
ou municípios. A parte interessada em um processo administrativo da SEAD-PI é, portanto, tipicamente
um servidor público, e não um cidadão na condição de usuário externo de serviço. A consequência
normativa dessa composição é discutida na seção 2.2 e delimita o alcance da afirmação jurídica que
o trabalho sustenta; a consequência metodológica é declarada aqui, no ponto em que o cenário é
caracterizado.

O procedimento administrativo estadual é regido pela Lei estadual n\textordmasculine{} 6.782, de 28
de março de 2016 \cite{state2016procedurelaw}, que disciplina o processo administrativo no âmbito
da Administração Pública do Estado do Piauí. É essa norma que estabelece as garantias
procedimentais --- forma, prazo, motivação, publicidade --- cuja verificação empírica pressupõe
justamente o registro que este trabalho produz.

O suporte material desse procedimento é o Sistema Eletrônico de Informações (SEI). O SEI é a
solução oficial do Governo Federal para produção e gestão de documentos e processos administrativos
eletrônicos, desenvolvida pelo Tribunal Regional Federal da 4\textordfeminine{} Região e cedida
gratuitamente a instituições públicas desde 2013 \cite{brasil2024sei}, no âmbito da infraestrutura
do Processo Eletrônico Nacional \cite{brasil2024pen}; o marco normativo que tornou obrigatório o
processo administrativo eletrônico na esfera federal, e que deu origem institucional a essa
infraestrutura, é o Decreto n\textordmasculine{} 8.539, de 8 de outubro de 2015
\cite{brasil2015decreto8539}. A cessão gratuita a órgãos estaduais e municipais fez do SEI o
substrato de fato do processo administrativo eletrônico brasileiro, e é nele que a SEAD-PI gera, instrui e conclui seus processos.

Três achados da literatura sobre o SEI sustentam premissas deste trabalho. Primeiro, o SEI já foi objeto de
revisão sistemática, que consolida os ganhos de eficiência e sustentabilidade relatados e, mais
relevante aqui, os problemas persistentes de interoperabilidade, segurança da informação e
preservação digital \cite{soriano2025sei}. Segundo, estudo de caso em universidade federal conclui
por ganho de eficiência na tramitação, mas registra que o SEI \emph{não} é um sistema informatizado
de gestão arquivística de documentos, do que decorrem risco de preservação e inconsistências de
classificação \cite{silva2020gestao}; é exatamente essa fragilidade que torna insuficiente confiar
no tipo de processo declarado no momento da autuação. Terceiro, a análise jurídica do processo
administrativo eletrônico brasileiro, com foco explícito no SEI, sustenta que a digitalização não
resolveu os problemas burocráticos anteriores e pode consistir em mera transposição de burocracia
para o meio eletrônico \cite{baptista2023processo} --- o que reforça a necessidade de tornar o
comportamento efetivo do procedimento inspecionável antes de automatizá-lo.

Há ainda um recorte útil do estado da arte: a literatura empírica sobre o SEI mediu sobretudo
\emph{percepção de usuário} --- adoção, vantagem relativa percebida, satisfação
\cite{fulton2018adocao} ---, e não o \emph{comportamento} do processo administrativo consignado nos registros de atividades. Trabalhos brasileiros de mineração de processos sobre registros de atividades extraídos do próprio SEI
existem, mas são raros \cite{limaesilva2026analise}. É nessa lacuna que este estudo se situa.

Três traços da implantação do SEI na Secretaria determinam o que o registro pode e não pode
mostrar, e condicionam as decisões de método das seções seguintes.

O primeiro é a \textbf{estrutura de unidades}. No SEI, \emph{unidade} é a entidade organizacional
que recebe, movimenta e conclui um processo: na SEAD-PI, correspondem a setores, diretorias e
gerências, dispostos em hierarquia e segmentados por objetivo e área de atuação. É essa entidade,
e não a pessoa, que figura no registro de tramitação --- razão pela qual o alfabeto de roteamento
da seção~\ref{sec:alfabetos} é composto de unidades, e não de servidores.

O segundo é que \textbf{um mesmo usuário pode atuar em nome de mais de uma unidade}. O vínculo
entre pessoa e unidade é de múltipla lotação, e não de correspondência um a um. Disso decorre que
uma transição entre unidades no registro não implica, necessariamente, transferência entre
pessoas distintas: o que o traço descreve é o percurso organizacional do processo, não a cadeia
de agentes que o conduziram.

O terceiro tem consequência direta sobre a representação adotada. Um protocolo do SEI ---
comumente chamado de processo --- guarda suas ações em ordem cronológica, mas pode correr
\emph{em paralelo} em mais de uma unidade: aberto um processo, ele pode ser enviado a várias
unidades ao mesmo tempo, cada uma praticando atos sobre ele sem que a outra conclua os seus. O
percurso de um processo não é, portanto, uma fila estritamente linear de unidades, e sim um
conjunto de trechos que pode ramificar e reconvergir. O alfabeto de roteamento da
seção~\ref{sec:alfabetos} é uma \emph{linearização} desse percurso pela ordem cronológica dos
eventos, e não uma reprodução de sua estrutura de concorrência. É uma simplificação assumida:
duas unidades que atuam simultaneamente aparecem no traço como se uma tivesse sucedido a outra,
com a ordem decidida pelo instante de cada evento. A escolha mantém o traço na forma de sequência
que os modelos da seção~\ref{sec:caminhos} exigem, ao custo de não distinguir sucessão de
simultaneidade --- limitação retomada na seção 4.11.

\subsection{Corpus e recorte}\label{sec:corpus}

O corpus provém do \emph{data lakehouse} da SEAD-PI sobre o SEI institucional do governo do
Estado do Piauí, hospedado pela ETIPI. A leitura se dá sobre esse espelho analítico --- as tabelas
Iceberg \texttt{sei\_bronze.sei\_atividades} e \texttt{sei\_gold.processos\_metadados}, consultadas
via Trino ---, e não sobre a base transacional do sistema de origem. O acervo ali disponível abrange
mais de vinte e um milhões de atividades entre 2021 e 2026; o recorte de um único ano adotado neste
trabalho é, portanto, escolha de projeto, e não limite da fonte --- ponto retomado na seção 4.11.

O predicado de entrada é \emph{nascido no órgão}: conta-se o processo
cujo evento de geração foi executado por uma unidade da SEAD-PI dentro do ano, e não todo processo
que apenas passou por ela, nem o campo de autuação do SEI. Um processo administrativo entra na
análise se satisfaz esse predicado, se foi concluído dentro do mesmo ano e se possui documento
inicial recuperável. O funil completo importa mais do que a sua última etapa, e é por isso dado
na íntegra: \textbf{12.621} processos nasceram na Secretaria durante o ano; \textbf{5.445}
compõem o corpus carregado; e \textbf{5.243} carregam documento inicial legível e entram na análise.

Duas medidas de ``concluído'' convivem nas fontes e respondem a perguntas distintas; cada uma é
reportada em separado. A do grafo (seção~\ref{sec:ekg}) exige que nenhuma unidade mantenha o processo em
aberto ao fim da janela; a do \emph{lakehouse} é textual, e reconhece a conclusão pela descrição do
andamento. A segunda registra \textbf{7.920} conclusões dentro do ano, contra os \textbf{4.701} que
não concluíram --- números que respondem a uma pergunta diferente da do corpus carregado. O
artefato que produz este funil carrega a mesma advertência junto dos valores.

Os processos que não haviam sido concluídos até 31 de dezembro são excluídos por construção,
e não por estimativa, e são estruturalmente os longos. A seção~\ref{sec:vies} quantifica uma segunda
perda, distinta desta, ocorrida \emph{dentro} da população concluída; ela não limita a primeira.

O que o corpus contém condiciona o alcance do artefato. Como registrado na seção~\ref{sec:cenario}, a Secretaria responde por pessoal e
administração geral, e sua carga é predominantemente de procedimento administrativo interno. A
versão curta da consequência normativa é que o catálogo aqui derivado instancia o dever
decorrente da carta de serviços apenas para o subconjunto de procedimentos que sejam serviços
prestados a usuário externo --- e o corpus não é principalmente composto por eles.

Restringir a análise a um único ano foi, em parte, decisão de custo para um estudo inicial: cinco anos de registros de atividades multiplicariam a ingestão, a vetorização, as varreduras de parâmetros de
agrupamento e uma matriz de avaliação que já leva cerca de onze horas para ser executada. Mas a
restrição não é apenas econômica: a cláusula 2.1 do instrumento autorizador delimita o objeto aos
processos nascidos e movimentados na SEAD-PI no período de janeiro a dezembro de 2025 (Anexo A). A
janela do estudo e a janela da autorização coincidem por construção.

\subsubsection{Censura à direita, medida}\label{sec:vies}

Exigir que um processo administrativo seja gerado \emph{e} concluído dentro de um mesmo ano civil
trunca a amostra, e o efeito desse truncamento é medido. O tempo
disponível para concluir encolhe mês a mês: um processo gerado em janeiro dispõe de 365 dias; um gerado em dezembro, de 31. Os máximos observados acompanham esse teto quase exatamente --- 343
dias contra 365 disponíveis em janeiro, 29 contra 31 em dezembro ---, o que é a assinatura da
censura, e não de quanto esses procedimentos de fato duram. A Tabela~\ref{tab:corpus-bias} dá o
quadro mês a mês.

%% GENERATED by `pipeline run results` - do not edit by hand.
\begin{table}[htbp]
\centering
\caption{Censura à direita por mês de abertura.}\label{tab:corpus-bias}
\footnotesize
\begin{tabular}{lrrrrr}
\toprule
mês & observados & dias disponíveis & vida mediana & vida máxima & excl.\ estimada \\
\midrule
2025-01 & 474 & 365 & 21 & 343 & 0,0\% \\
2025-02 & 619 & 334 & 14 & 290 & 0,2\% \\
2025-03 & 478 & 306 & 17 & 284 & 0,8\% \\
2025-04 & 470 & 275 & 25 & 244 & 2,1\% \\
2025-05 & 415 & 245 & 24 & 225 & 5,1\% \\
2025-06 & 621 & 214 & 18 & 204 & 21,7\% \\
2025-07 & 653 & 184 & 12 & 166 & 24,1\% \\
2025-08 & 398 & 153 & 18 & 129 & 26,2\% \\
2025-09 & 539 & 122 & 10 & 105 & 29,3\% \\
2025-10 & 381 & 92 & 10 & 75 & 31,9\% \\
2025-11 & 268 & 61 & 5 & 42 & 35,9\% \\
2025-12 & 129 & 31 & 5 & 29 & 44,5\% \\
\bottomrule
\end{tabular}
\par\vspace{2pt}\footnotesize\raggedright Observados 5.445; verdadeiro estimado 6.658; ausentes estimados \textbf{1.213} (18,2\%). Limite inferior: o próprio mês de referência está truncado.
\end{table}


Tomando o mês menos truncado como distribuição de referência, o número de processos que teriam sido
concluídos dentro de uma janela de comprimento pleno é estimado em 6.658, contra os 5.445
observados; estima-se, portanto, que 18,2\% da população \emph{concluída} esteja ausente, e que
sejam sistematicamente os processos longos. A cifra é um déficit \emph{dentro} dos
processos que concluem, estimado a partir da censura mês a mês. Nada diz sobre os 7.176 processos
que jamais concluíram dentro da janela, que são excluídos de saída (seção~\ref{sec:corpus}). É também um limite
inferior dentro do seu próprio alcance, já que o mês de referência é ele próprio truncado. Três
consequências vinculam o restante do trabalho e são retomadas na seção 4.11: o corpus não é ``processos gerados em 2025'', e sim ``processos gerados em 2025 que também concluíram rápido o
bastante''; nenhuma estatística de duração ou de eficiência derivada dele é admissível; e
procedimentos longos, que atravessam várias unidades, estão sub-representados.

\subsection{Grafo de conhecimento de eventos}\label{sec:ekg}

Os dados de eventos são modelados segundo a prática consolidada em mineração de processos
\cite{vanderaalst2016processmining}. \textbf{O que entra.} A etapa de extração lê, do espelho
analítico descrito na seção~\ref{sec:corpus}, três conjuntos: as atividades registradas pelo sistema
de origem, o cadastro de unidades organizacionais e os metadados de cada processo administrativo. O
recorte ali definido é aplicado já na leitura --- órgão, ano civil e condição de conclusão ---, de
modo que nada além dele é materializado.

\textbf{O que o grafo materializa.} A carga procede em cinco passos encadeados, em ordem fixa.
Primeiro os nós-dimensão --- unidades organizacionais, usuários e tipos de
procedimento declarados --- são pré-criados, e assim as cargas seguintes se fundem a nós
estáveis em vez de duplicá-los. Em seguida as atividades são anexadas: cada uma é ligada ao seu
processo administrativo, à unidade que a executou, ao usuário que a praticou e ao seu tipo de ação.
Os nós de processo nascem do evento de geração do próprio processo, o que fixa uma raiz por caso
antes de qualquer atividade ser ligada a ela, e os documentos produzidos ou referenciados são
ligados ao processo que os contém. O terceiro passo é o que faz do registro um processo: as
atividades de um mesmo caso são encadeadas por uma relação de precedência direta, na ordem
cronológica do registro \cite{esser2021ekg}. Sobre essa cadeia derivam-se, por fim, as grandezas de
posse --- quanto tempo cada unidade reteve cada processo e quantas vezes o recebeu de volta, que é
onde o retrabalho se torna visível --- e a situação de cada caso ao fim da janela. É neste último
passo que se decide a conclusão pela ausência de vínculo em aberto em qualquer unidade, e não pelo
registro textual de um evento de encerramento: são as duas medidas incompatíveis da
seção~\ref{sec:corpus}, e o grafo adota a primeira.

\textbf{O que sai.} O grafo é a fonte única das etapas seguintes: dele saem, por travessia da
cadeia de precedência, os traços nos três alfabetos da seção~\ref{sec:alfabetos}, e dele saem também
os blocos de atributos que descrevem o percurso --- unidades visitadas, transições entre unidades e
histograma de tipos de ação --- da seção~\ref{sec:atributos}. Derivar os traços de um grafo, em vez
de um registro de atividades plano, é o que mantém simultaneamente disponíveis as perspectivas de
unidade e de ação, das quais os três alfabetos dependem. Uma etapa dedicada congela ainda esse
substrato de traços em disco antes da avaliação, de modo que a matriz inteira e qualquer
reanálise posterior corram sem banco de dados, sobre exatamente os mesmos traços. Ainda assim,
adotou-se uma noção de caso
--- o processo administrativo ---, de modo que as preocupações com divergência e convergência que
motivam a mineração de processos centrada em objetos permanecem aplicáveis
\cite{vanderaalst2019ocpm}.

\subsection{O documento inicial: leitura, supressão de dados pessoais e destilação}\label{sec:documento}

A anonimização é requisito do instrumento que autoriza o estudo, e o requisito alcança inclusive a
\emph{ordem} das etapas. A cláusula 5.1 do Termo de
Ciência e Responsabilidade \cite{seadpi2026termo} dispõe:

\begin{citacao}
5.1. A metodologia incluirá etapa de pseudonimização e/ou anonimização dos dados pessoais, de
forma automatizada sempre que identificáveis, antes das fases de clusterização e vetorização.
\end{citacao}

O Anexo A reproduz a cláusula em seu contexto, ao lado das demais que vinculam este trabalho. Dela
decorrem duas obrigações distintas: que exista uma etapa automatizada de supressão de dados
pessoais, e que ela \emph{preceda} a vetorização e o agrupamento. Ambas são condições de licitude
do tratamento, tal como pactuadas.

O instrumento invoca, ao fazê-lo, o art. 12 da Lei Geral de Proteção de Dados
\cite{brasil2018lgpd}, e a leitura desse dispositivo importa. O art. 12 é uma norma
\emph{delimitadora de escopo} --- dados anonimizados ficam fora do alcance da lei, salvo quando o
processo puder ser revertido por meios razoáveis --- e não um dever autônomo de anonimizar. O dever,
aqui, é contratual; o que o art. 12 fornece é o padrão contra o qual o resultado é julgado, e a
consequência de não atendê-lo, que é a devolução dos dados ao alcance integral da lei.

O custo analítico da transformação é medido: o \emph{pipeline} é executado duas vezes --- uma sobre
o texto bruto, outra sobre a variante anonimizada --- e os dois braços são comparados processo a
processo. A seção 4.6 reporta o resultado dessa comparação.

O documento inicial é primeiro convertido a texto plano --- extração direta quando o arquivo é
estruturado, leitura por modelo de visão computacional executado localmente quando não é ---, e é
sobre esse texto que a supressão incide. A detecção de dados pessoais combina duas camadas. A
primeira é de regras, inspirada em mecanismos baseados em reconhecedores declarativos
\cite{microsoftpresidio}: identificadores fiscais de pessoa física e jurídica são reconhecidos nas
formas com e sem pontuação e \emph{validados pelos próprios dígitos verificadores}, o que impede
que números de protocolo e de documento sejam suprimidos por semelhança de formato; endereços de
correio eletrônico seguem o padrão da especificação; e a matrícula funcional só é reconhecida em
contexto, isto é, precedida do termo que a designa e com extensão mínima. A segunda camada é um
extrator neural de entidades nomeadas com esquema aberto \cite{zaratiana2025gliner2}, cujas
extensões detectadas são mapeadas sobre o mesmo conjunto de rótulos tipados. O mapeamento é largo:
qualquer chave livre de identidade ou de finança que o extrator emita recai em um rótulo genérico.

A substituição preserva a integridade referencial: o marcador é derivado do resumo criptográfico do
par rótulo-entidade em forma canônica, de modo que a mesma pessoa recebe o mesmo marcador em todos
os documentos em que apareça, ao passo que a forma de superfície desaparece. As ocorrências são
substituídas da extensão mais longa para a mais curta, com as guardas de fronteira de cada detector
--- dígito para os rótulos numéricos, caractere de palavra para os nominais ---, porque uma guarda
única deixaria escapar em silêncio as ocorrências adjacentes a dígitos. Um modo de verificação
reexecuta a camada de regras \emph{sobre a saída} e reprova a execução se restar qualquer
ocorrência.

Suprimir nomes, porém, não anonimiza por si só registros de atividades: um traço é, ele próprio, uma
impressão digital comportamental \cite{mannhardt2019privacy, elkoumy2022privacy}, e uma sequência
datada de movimentações entre unidades organizacionais nominadas permanece um registro do
comportamento dos servidores que trataram do caso.

\textbf{O que ela cobre.} Identificadores diretos e indiretos presentes no texto livre do documento
inicial --- nomes de pessoas, identificadores nacionais, endereços, dados de contato, números de
matrícula --- são detectados e substituídos por marcadores tipados antes que qualquer vetorização
seja computada. \textbf{O que ela não cobre.} O traço permanece intocado: a sequência de unidades
organizacionais, os tipos de ação e seus carimbos de tempo são todos preservados, porque são o
objeto de estudo. Um leitor com acesso ao sistema de origem poderia reidentificar um processo
administrativo apenas pela sua sequência de roteamento. Os dois braços diferem, portanto, apenas no
canal \emph{textual}, e é dentro desse escopo que a comparação da seção 4.6 deve ser lida.
``Anonimizado'', ao longo do texto, designa essa transformação específica.

Do texto assim tratado o método destila, por fim, duas representações mais curtas. Um modelo de
linguagem de pesos abertos produz, para cada documento inicial, um \textbf{tópico} --- o assunto em
poucas palavras --- e um \textbf{resumo} de uma frase, ambos conservados ao lado do texto. A
destilação existe porque o vetor do texto integral carrega, além do assunto, o formulário
administrativo em que ele vem embrulhado; tópico e resumo são um sinal de assunto com menos ruído,
e entram como bloco próprio na representação da seção~\ref{sec:atributos}. Como o tópico é derivado
do texto, ele difere entre os dois braços, e por isso é recomputado a cada troca de braço em vez de
reaproveitado --- reaproveitá-lo faria o braço anonimizado herdar o assunto extraído do texto
bruto, e a comparação entre braços deixaria de medir o que declara medir.

\subsection{Vocabulário controlado de palavras-chave}\label{sec:keywords}

Do mesmo texto de que sai a vetorização sai também um objeto de natureza distinta, e é o único de
todo o método que uma pessoa consegue ler: o vocabulário controlado de palavras-chave. Ele é a
ponte entre o vetor que agrupa e a prosa que nomeia, e sem ele a nomeação da
seção~\ref{sec:nomeacao} regride a amostras de documentos. A pergunta que ele responde é que termos
o próprio acervo usa para se descrever, sem que a lista seja escrita à mão por quem analisa; por
isso o vocabulário é \emph{crescente}, e não fixado a priori: o modelo de linguagem propõe termos
documento a documento, e cada termo proposto ou se funde a um termo já existente ou entra como
novo. O vocabulário resultante é, por construção, idêntico nos dois braços.

Um termo entra no vocabulário por três filtros encadeados. Primeiro, a \textbf{normalização}:
redução a minúsculas, remoção de acentos, descarte de pontuação, colapso de espaços, caminhos de
unidade reduzidos à sigla final e rejeição de expressões meramente temporais, como um ano isolado
ou o exercício a que ele se refere. Em seguida, o \textbf{descarte de dois conjuntos} --- as
palavras vazias da língua e os termos que descrevem a própria administração: o nome do órgão, as
suas unidades, o vocabulário de tramitação. Estes aparecem em quase todo documento e por isso não
separam coisa alguma; mantê-los daria ao bloco um eixo constante. Por fim, a \textbf{fusão de
variantes}: dois termos cuja similaridade de Dice sobre bigramas de caracteres alcance o limiar de
0,9 são tratados como um só, e o primeiro a ser observado torna-se a forma canônica. É esse limiar
que decide se duas grafias próximas de um mesmo assunto são um termo ou dois. Uma segunda passagem
reanota os documentos processados enquanto o vocabulário ainda se construía, de modo que nenhum
documento fique anotado contra um vocabulário parcial. Entra no bloco de atributos apenas o termo
que ocorra em ao menos três processos administrativos, e é essa exigência que fixa a largura do
bloco.

Três propriedades decorrem desse procedimento. A similaridade é
\textbf{léxica}, e erra em duas direções previsíveis: funde variações morfológicas de um mesmo
termo, o que se pretende, mas não funde sinônimos de grafia distinta e pode fundir termos curtos
que apenas se parecem. A escolha da forma canônica \textbf{depende da ordem de chegada} dos
documentos, o que torna a etapa não determinística entre execuções; por isso o vocabulário é
persistido como artefato e reutilizado, e é ele --- e não o procedimento --- que garante a
reprodutibilidade dos blocos que dele dependem. E a etapa precisa ser executada \textbf{uma única
vez, antes} da construção de ambos os braços: executada depois, o bloco de palavras-chave chega
vazio à vetorização, o agrupamento prossegue sobre um bloco nulo e nenhuma verificação acusa a
falta --- razão pela qual a ordem entre esta etapa e a construção dos braços é tratada como
invariante do procedimento, e não como conveniência de execução.

\subsection{Blocos de atributos}\label{sec:atributos}

A afirmação substantiva do método está nesta etapa: forma procedimental e conteúdo documental são
ambos sinal, e nenhum dos dois isoladamente identifica um tipo de processo. Cada processo
administrativo é, portanto, representado por seis blocos que fundem o que o caso \emph{diz} com o
modo como ele \emph{se moveu}:

\begin{itemize}
  \item o documento inicial vetorizado por um codificador multilíngue de sentenças
    \cite{chen2024bgem3} (1.024 dimensões);
  \item um tópico e um resumo desse documento produzidos por modelo de linguagem, vetorizados do
    mesmo modo (1.024);
  \item TF--IDF sobre as unidades organizacionais que o processo visita (500) e sobre as transições
    entre unidades (1.000);
  \item um histograma sobre os tipos de ação (48);
  \item TF--IDF sobre um vocabulário controlado de palavras-chave extraído do documento inicial
    (517).
\end{itemize}

Os dois blocos de conteúdo usam \textbf{apenas o documento inicial}: os
documentos posteriores são majoritariamente formulário de tramitação --- encaminhamentos, despachos
de rotina --- que se parecem entre milhares de processos sem relação alguma e diluem o centroide,
ao passo que o documento de abertura é aquilo que o caso declara ser antes de o procedimento ter
feito qualquer coisa com ele. Um processo sem documento legível recebe blocos de conteúdo nulos e
ainda assim é agrupado pela sua assinatura estrutural. Nos blocos estruturais, o primeiro registra a
presença de cada unidade no percurso, e não quantas vezes ela aparece; o segundo toma os pares
consecutivos de unidades depois de colapsadas as repetições, de modo a descrever transição e não
permanência; e ambos são ajustados sobre o próprio lote, o que faz a importância de um termo
refletir o conjunto agrupado.

Cada bloco é normalizado em $L^2$ por linha antes da concatenação e todos os seis carregam peso
1,0, e assim o esquema pondera igualmente os blocos, independentemente de sua largura: o
histograma de 48 tipos de ação contribui com tanta massa quanto a vetorização de 1.024 dimensões do
documento. Cada peso é exposto como parâmetro, e é isso que torna a afirmação de fusão testável
pela remoção controlada de blocos da seção~\ref{sec:evidencia}: um bloco pode ser zerado sem que
nada mais mude.

O vetor concatenado, com mais de três mil dimensões, é então reduzido a \textbf{50} por UMAP
\cite{mcinnes2018umap} sob distância de cosseno --- a métrica que corresponde ao espaço de
vetorizações normalizadas ---, com 15 vizinhos e semente fixa. A redução reproduz-se bit a bit na
mesma máquina, mas não entre processadores e versões de biblioteca distintos; a etapa declara essa
condição no artefato que escreve, junto da semente, de modo que o leitor saiba de antemão quais
grandezas reproduzem exatamente e quais reproduzem apenas em distribuição.

Tanto o modelo de linguagem que produz tópicos e palavras-chave quanto o modelo de vetorização são executados, a partir de pesos abertos, na infraestrutura computacional do próprio Estado: a inferência ocorre no \textbf{Mandu}, o \emph{cluster} de alto desempenho com GPUs NVIDIA do Governo do Estado do Piauí, que expõe uma API compatível com o padrão OpenAI. A compatibilidade de API é conveniência de integração, e não indício de externalização. Isso é requisito jurídico, e não preferência de engenharia: a cláusula 4.1 do instrumento autorizador exige que o tratamento permaneça dentro da infraestrutura do Estado, e a cláusula 4.5, alínea b, veda interfaces externas que impliquem saída de dados. Reapontar o endereço-base da API para um provedor externo violaria o termo.

\subsection{Agrupamento}\label{sec:agrupamento}

\subsubsection{Três famílias, e por que três}\label{sec:familias}

O agrupamento é conduzido com três famílias de algoritmo: baseada em \textbf{densidade}
\cite{campello2013hdbscan, mcinnes2017hdbscan}, baseada em \textbf{centroides}
(\textit{k}-médias) e baseada em \textbf{modularidade} \cite{traag2019leiden}. Executar três
famílias distingue uma propriedade dos dados do viés indutivo de um algoritmo particular. Uma estrutura que só aparece sob uma das três é um achado sobre o
algoritmo; uma que aparece sob as três é um achado sobre o acervo.

As três operam sobre a mesma matriz de 50 dimensões da seção~\ref{sec:atributos}, e os parâmetros
são varridos, e não fixados. A família de densidade dispensa que se declare o número de grupos:
ela constrói a hierarquia de densidade dos pontos e dela extrai a partição mais estável, com
tamanho mínimo de grupo igual a 20 e seleção por \emph{excess of mass}; uma variante mais grossa
funde grupos separados por menos que $\varepsilon = 0{,}5$. A métrica de cosseno descreve a \emph{redução}, ao passo que o agrupamento por
densidade corre em distância euclidiana sobre o espaço já reduzido --- prática usual, uma vez que a
projeção preserva vizinhança e não ângulo, mas que não deve ser lida como agrupamento por cosseno.
A família de modularidade opera sobre um grafo de $k$ vizinhos mais próximos ($k = 15$, cosseno),
com a resolução varrida entre 0,4 e 2,0 em passos de 0,2. A família de centroides varre $k$ de 20 a
320 em passos de 10, com dez reinicializações por valor e semente fixa, sob os dois critérios de
parada descritos a seguir. A varredura produz \textbf{sete} parametrizações por braço ---
\textbf{catorze} ao todo ---, detalhadas na seção~\ref{sec:desenho}.

\subsubsection{A escolha de $k$ e da resolução, varrida e reportada}\label{sec:selecaok}

O número de grupos é escolhido por varredura: uma etapa dedicada percorre a grade de $k$ e reporta,
lado a lado, o critério do cotovelo --- definido como a máxima distância perpendicular à corda que
liga os extremos da curva de inércia, com ambos os eixos normalizados ao intervalo unitário, e não
por leitura visual --- e a \emph{silhouette}. Ambos os valores são conservados como parametrizações
distintas, em vez de se eleger um deles \emph{a
priori}: os dois critérios otimizam objetivos diferentes --- compacidade relativa à separação, num
caso; ganho marginal de inércia, no outro.
Levar os dois adiante transforma a escolha de $k$ de premissa em objeto de avaliação. A resolução da
família de modularidade é varrida do mesmo modo, e pelo mesmo motivo, elegendo-se a que maximiza a
\emph{silhouette} de cosseno. Onde o cálculo da \emph{silhouette} sobre a matriz completa exceda o
orçamento de tempo fixado, ele passa a ser feito sobre uma subamostra de tamanho fixo e semente
fixa, e o escopo empregado é registrado no próprio artefato --- de modo que um valor de subamostra
nunca seja lido como se fosse de matriz completa. A varredura inteira é persistida como artefato
auditável, e não apenas o seu vencedor.

\subsubsection{O ruído da família de densidade é atribuído}\label{sec:ruido}

A família de densidade é a única das três que pode recusar-se a atribuir um processo administrativo
a qualquer grupo. Essa recusa é informativa, mas tem uma consequência que inviabiliza a comparação
entre famílias: os índices de validade e a acurácia a jusante passariam a ser calculados sobre
\emph{populações diferentes}, e uma família que avalia apenas os casos fáceis é premiada por isso.

Por omissão, portanto, cada ponto recusado é atribuído ao grupo cujo centróide lhe é mais próximo,
no mesmo espaço em que as métricas são calculadas. Todas as famílias passam a ser avaliadas sobre os
mesmos casos, e a comparação entre elas torna-se legítima. A escolha tem um custo: um processo que
não pertence a família alguma passa a ser arquivado na mais
próxima, o que é a troca correta para a avaliação e a errada para um catálogo que um servidor
folheia. O que a densidade \emph{teria} recusado não se perde --- fica registrado, por agrupamento,
como a fração recusada antes da atribuição, e é essa grandeza, e não a taxa de ruído final, que
responde por abstenção em qualquer análise. A alternativa de reclassificar os recusados em
subgrupos próprios permanece implementada e recuperável por parâmetro.

\subsubsection{Partições-controle embaralhadas, uma por agrupamento}\label{sec:controles}

Para cada um dos catorze agrupamentos reais é construída uma \textbf{partição-controle}: os rótulos
são permutados entre os processos administrativos, preservando exatamente o número de grupos, a
distribuição de tamanhos e a população avaliada, e destruindo apenas a pertença. É contra ela, e não
apenas contra um modelo global, que o ganho é medido (seção~\ref{sec:estimando}).

O controle é construído \emph{por agrupamento}, e não como um único nulo compartilhado, por uma
razão substantiva: a forma de uma partição --- quantos grupos, de que tamanhos --- desloca a
acurácia por si só, independentemente de a pertença significar alguma coisa. Um controle derivado de
\emph{outro} agrupamento mediria diferença de granularidade em vez de diferença de agrupamento. Por
isso a partição-controle é artefato da fase de agrupamento, e não da fase de avaliação.

\subsection{Alfabetos de traço}\label{sec:alfabetos}

O traço de um processo administrativo é renderizado em \textbf{três} alfabetos, e a distinção
delimita o que cada número adiante pode afirmar. O alfabeto de \textbf{ação por unidade} emparelha
cada ação registrada com a unidade que a executou --- a sigla completa da unidade, sem qualquer
agregação por diretoria ---, e é o estado não projetado; os outros dois são projeções dele, obtidas
tomando um dos dois eixos e colapsando repetições consecutivas. O alfabeto de \textbf{roteamento}
conserva apenas
a sequência de unidades organizacionais que o processo visita, abstraindo o que ocorreu dentro de
cada uma --- uma perspectiva organizacional consolidada \cite{song2008organizational}. O alfabeto
de \textbf{procedimento} faz o inverso: conserva a sequência de ações e abstrai o organograma, o que
o torna a mais transferível das três representações, por não depender dos nomes de unidade desta
administração.

Colapsar repetições define o roteamento: um caso que permanece
quatro eventos na mesma unidade fez \emph{uma} visita, não quatro, e uma unidade que reaparece de
forma não consecutiva é um retorno de verdade. A permanência não desaparece do método por isso ---
a contagem de eventos por visita é conservada ao lado do traço ---, apenas sai do alfabeto. Cada
traço recebe ainda um estado sentinela de encerramento, sem o qual um gerador não teria como
aprender \emph{quando} um caso termina, e pararia no primeiro estado de conclusão que encontrasse.
Três alfabetos, e não um, porque respondem a perguntas administrativas distintas --- o que se faz,
onde se faz e por onde o caso anda --- e um método que sirva bem a uma delas pode falhar nas
outras.

\subsubsection*{Cada alfabeto é avaliado dentro de si}

Cada alfabeto recebe a sua própria matriz completa de avaliação, e o melhor modelo é eleito
\emph{dentro} de cada um. Nenhum número atravessa alfabetos em nenhum ponto deste trabalho. A
decisão assenta em três razões independentes:

\begin{enumerate}
  \item \textbf{As populações diferem.} As duas projeções colapsam repetições consecutivas, de modo
    que um processo administrativo tratado inteiramente dentro de uma única unidade reduz-se a um
    só estado de roteamento e é descartado por não ter transição alguma a prever. O alfabeto de
    roteamento avalia, portanto, os processos multiunidade --- não o corpus.
  \item \textbf{Os comprimentos diferem.} Toda medida sensível ao comprimento do traço, e as
    medidas de sufixo o são, muda de escala entre alfabetos sem que nada tenha mudado no acervo.
  \item \textbf{Os vocabulários diferem por ordens de grandeza.} Uma acurácia de próximo passo
    maior sobre um alfabeto menor pode significar apenas menos candidatos por passo. Comparar
    acurácias sobre vocabulários de tamanhos distintos compara a dificuldade da tarefa, e não a
    qualidade do método.
\end{enumerate}

Um resultado de roteamento é uma afirmação sobre percurso entre unidades; um resultado de
procedimento é uma afirmação sobre o que se faz.

Um exemplo torna a diferença concreta. Um processo do corpus é gerado em uma unidade, tem
documentos gerados e assinados ali, segue para uma segunda unidade, recebe novo documento, retorna
à unidade de origem e é concluído. A Figura~\ref{fig:alfabetos} mostra como esse mesmo processo é
renderizado nos três alfabetos; as unidades reais foram substituídas por rótulos genéricos, e a
unidade de origem reaparece como \texttt{UNIDADE-A} nos dois pontos em que o processo passa por ela.

\begin{figure}[htbp]
\centering
\fbox{\begin{minipage}{0.93\textwidth}
\vspace{2pt}
\raggedright\footnotesize
\textbf{Alfabeto de ação por unidade} --- excerto dos dezoito estados:

\vspace{2pt}
\texttt{GERACAO-PROCEDIMENTO@UNIDADE-A} · \texttt{GERACAO-DOCUMENTO@UNIDADE-A} ·
\texttt{ASSINATURA-DOCUMENTO@UNIDADE-A} · \texttt{CONCLUSAO-AUTOMATICA@UNIDADE-A} ·
\texttt{PROCESSO-REMETIDO@UNIDADE-B} · \texttt{PROCESSO-RECEBIDO@UNIDADE-B} ·
\texttt{GERACAO-DOCUMENTO@UNIDADE-B} · \ldots ·
\texttt{CONCLUSAO-PROCESSO@UNIDADE-A}

\vspace{6pt}
\textbf{Alfabeto de roteamento} --- os mesmos dezoito estados, colapsados em três:

\vspace{2pt}
\texttt{UNIDADE-A} $\rightarrow$ \texttt{UNIDADE-B} $\rightarrow$ \texttt{UNIDADE-A}

\vspace{6pt}
\textbf{Alfabeto de procedimento} --- os mesmos dezoito estados, abstraídos do organograma:

\vspace{2pt}
\texttt{GERACAO-PROCEDIMENTO} $\rightarrow$ \texttt{GERACAO-DOCUMENTO} $\rightarrow$
\texttt{ASSINATURA-DOCUMENTO} $\rightarrow$ \texttt{PROCESSO-REMETIDO} $\rightarrow$
\texttt{PROCESSO-RECEBIDO} $\rightarrow$ \texttt{GERACAO-DOCUMENTO} $\rightarrow$ \ldots
$\rightarrow$ \texttt{CONCLUSAO-PROCESSO}
\vspace{3pt}
\end{minipage}}
\caption{Um mesmo processo administrativo renderizado nos três alfabetos de traço. O alfabeto de
ação por unidade conserva o que se fez e onde; o de roteamento conserva apenas por onde o processo
passou; o de procedimento conserva apenas o que se fez. Os dois últimos colapsam repetições
consecutivas, e é por isso que avaliam populações distintas.}
\label{fig:alfabetos}
\end{figure}

\noindent O retorno à unidade de origem permanece visível nos alfabetos que conservam a unidade --- é essa repetição
que os modelos de caminho típico da seção~\ref{sec:caminhos} precisam aprender.

A convenção fixada na seção~\ref{sec:caracterizacao} vale aqui sem exceção: o alfabeto é o de
\emph{roteamento}; \emph{tramitação} fica reservada ao movimento administrativo que ele representa.

A motivação é que a unidade carrega significado que a ação não carrega. Saber que um processo
administrativo foi de um protocolo para uma unidade de parecer jurídico e desta para uma unidade
orçamentária descreve o procedimento; saber que um documento foi gerado, assinado e concluído
descreve convenções do sistema, comuns a quase todo procedimento.

\subsection{Modelos de caminho típico e decodificação restrita}\label{sec:caminhos}

Para cada grupo são ajustados \textbf{seis} geradores, que se distribuem por três modos distintos de
estimar o próximo estado. Dois são \textbf{contadores locais}: uma cadeia de Markov de terceira
ordem e um grafo de precedência direta (DFG, rótulo \texttt{dfg} nas tabelas). Três são
\textbf{sequenciais neurais}: GRU \cite{cho2014gru}, LSTM \cite{hochreiter1997lstm} e Transformer
\cite{vaswani2017attention}. O sexto é \textbf{convolucional sobre grafo} (GNN, rótulo \texttt{gnn}
nas tabelas) \cite{kipf2017gcn}, e generaliza por uma via que nenhum dos outros cinco usa.

As três famílias neurais são as três linhagens que a literatura de predição de próxima atividade
consolidou --- a recorrente \cite{tax2017lstm}, a recorrente com portas e a baseada em atenção
\cite{bukhsh2021processtransformer} ---, e as três são comparadas porque
sob protocolo comum nenhuma arquitetura domina em todos os registros de atividades
\cite{ramamaneiro2023dlppm}. Qual delas vence é, portanto, pergunta empírica endereçada ao corpus, e
não premissa do método. As três compartilham a mesma configuração, fixada uma única vez e idêntica
em todos os grupos e em todas as células de avaliação, sem ajuste por grupo: vetorização de estados
em 32 dimensões, estado oculto de 64, janela de contexto de 8 estados, \emph{dropout} de 0,2 e 30
épocas de Adam com taxa de aprendizado 0,01, decaimento de peso $10^{-4}$, lotes de 2.048 e semente
fixa; a GRU e a LSTM com uma camada recorrente, o Transformer com duas camadas de codificação e
quatro cabeças de atenção. Fixar a configuração de antemão é o que impede que a comparação entre
famílias meça esforço de ajuste em vez de arquitetura. Grupos com menos de 30 traços não são
avaliados. Rolar um
modelo sequencial treinado adiante para gerar um traço inteiro, em vez de pontuar predições de um
passo, é o arranjo de \citeonline{camargo2019lstm}, e o realismo da sequência gerada é o que as
medidas de sustentação de caminho, definidas na seção~\ref{sec:protocolo}, existem para aferir.

O gerador de
Markov é de terceira ordem com retrocesso: onde um contexto de terceira ordem não foi observado, ele
recua para segunda e depois para primeira ordem, o que é o que o mantém utilizável em grupos
próximos ao piso de 30 traços. Emprega-se um grafo de precedência direta como um dos aprendizes de
caminho --- de primeira ordem, restrito às transições do conjunto de treino que alcancem um piso de
cobertura de 0{,}05 --- a transição deve ocorrer em ao menos 5\% dos traços de treino do grupo ---,
registrando-se que essa representação tem patologias documentadas enquanto modelo de
processo \cite{vanderaalst2019dfg}.

O gerador convolucional entra por uma razão que os outros cinco não cobrem. Markov e precedência
direta são estimadores estritamente locais: a previsão para um estado usa apenas as transições
observadas a partir dele, de modo que um estado visto três vezes na partição tem uma estimativa de
três amostras e nada a corrige. Sobre o grafo de transições do conjunto de treino --- nós são
estados, arestas são transições contadas, a adjacência simetrizada e acrescida de laços próprios,
normalizada de forma simétrica --- duas rodadas de propagação fazem a previsão de um estado
incorporar a de seus vizinhos, de modo que um estado raro toma emprestado de estados
estruturalmente semelhantes. É a única família do conjunto que generaliza \emph{pelo grafo}, e não
pela sequência, e ela paga por isso condicionando apenas no último estado, sem janela de contexto:
ganha vizinhança e perde memória. Treina-se por entropia cruzada sobre os pares observados, com
Adam, estado oculto de 64, 200 épocas, taxa de aprendizado 0,01, decaimento de peso $10^{-4}$ e
semente fixa. A reprodutibilidade bit a bit exige ainda uma providência que não é óbvia: a
paralelização de álgebra linear é reduzida a uma única linha de execução durante o ajuste, porque
somas em ordem variável divergem o bastante para trocar a ordem das melhores predições. Herda dos
contadores todo o resto --- predição, distribuição do próximo estado, geração, feixe, amostragem e
decodificação restrita ---, o que mantém a comparação sobre a \emph{estimativa}.

O gerador de precedência direta compartilha estrutura com uma das linhas de base. Sob predição de
primeira escolha, ele seleciona o sucessor observado mais frequente, que é exatamente a linha de
base do sucessor intragrupo descrita na seção~\ref{sec:protocolo}; os dois são o mesmo estimador e
pontuam de forma idêntica. Há, portanto, cinco preditores treinados distintos mais uma linha de
base que também
comparece ao conjunto de geradores, e o máximo por grupo sobre ``seis famílias'' é tomado sobre um
conjunto que contém o seu próprio piso. Se essa sobreposição altera as cifras de melhor modelo
selecionado é questão empírica, respondida na seção~\ref{sec:res-sinal}.

A geração é decodificada sob uma máquina de estados que codifica as regras de movimentação do sistema de
processo eletrônico. A máquina de estados garante uma propriedade de segurança, e não uma condição de aceitação. A cada passo,
o decodificador admite apenas os estados que a máquina de estados permite, e assim um caminho produzido
não pode receber um processo que nunca lhe foi enviado, nem concluir uma raia que nunca foi aberta.
Ele não garante que todo envio seja correspondido antes de o caminho terminar: a terminação é
irrestrita, de modo que um caminho pode encerrar com uma transferência pendente. A restrição também
vincula apenas os cinco tipos de ação que participam das regras de movimentação, permanecendo
irrestritos os quarenta e três restantes. A seção 4.5 reporta o que a garantia, assim delimitada,
efetivamente compra.

\subsection{Nomeação dos tipos}\label{sec:nomeacao}

Uma entrada de catálogo precisa de nome, e o nome é a parte que o leitor vê primeiro. Os nomes aqui
são gerados, e não redigidos, e a nomeação corre em \textbf{duas passagens}. Nomear cada grupo
isoladamente a partir das suas
palavras-chave mais \emph{frequentes} garante nomes parecidos para grupos parecidos, porque numa
família de procedimentos vizinhos os termos mais frequentes são justamente os \emph{partilhados},
e o modelo, que nunca vê os concorrentes, não tem como diferenciá-los.

Na primeira passagem, portanto, cada grupo é apresentado ao modelo de linguagem executado
localmente por \emph{dois} conjuntos de palavras-chave em vez de um: as mais frequentes, que dão o
tema, e as \textbf{distintivas} --- as sobre-representadas naquele grupo em relação ao resto do
acervo, ordenadas pela razão entre a probabilidade do termo dentro e fora do grupo, com suavização
para os termos raros. São estas últimas que separam, ao passo que as primeiras unem. Junto delas
seguem uma amostra dos processos administrativos mais centrais --- o tópico e o resumo de cada
documento inicial ---, as unidades organizacionais visitadas e os tipos de ação; e o modelo devolve
um nome em português, uma descrição factual de uma frase e três marcadores
distintivos. O modelo é instruído a ancorar o nome no conteúdo e no fluxo do próprio grupo, a
preferir um assunto concreto a um rótulo genérico e a \emph{não} utilizar a classificação
declarada, para que o tipo emergente não seja silenciosamente rerrotulado com a taxonomia da qual o
estudo declina depender --- usá-la tornaria circular a concordância contra o tipo declarado, que é
justamente o instrumento externo de avaliação da partição.

A segunda passagem renomeia \emph{apenas} os grupos cujos nomes colidiram, e agora com as rivais à
vista: cada um recebe o nome, o tamanho e os termos distintivos dos seus homônimos, e a instrução
explícita de produzir um nome que os distinga. Nomes cujas vetorizações permaneçam quase duplicatas
são fundidos por agrupamento aglomerativo sobre a similaridade de cosseno entre as vetorizações dos
nomes; a fusão é uma etapa do procedimento. A nomeação é executada para \emph{todas} as partições, e não apenas para aquela que se
julgue vencedora, uma vez que o critério de seleção da seção~\ref{sec:desenho} só é aplicado depois
--- e uma partição sem nomes renderizaria um catálogo vazio.

Duas limitações se seguem. Os nomes são saída de máquina
não validada sobre matéria administrativa, o que constitui artefato distinto da própria partição e
carrega risco distinto; a seção 4.11 enuncia esse ponto ao lado da ausência de avaliação por
especialistas dos tipos obtidos. A nomeação quase duplicada foi observada durante o
desenvolvimento; a etapa de fusão a mitiga.

\subsection{Protocolo de avaliação e os três controles}\label{sec:protocolo}

\subsubsection*{A unidade de avaliação: a célula}

A unidade elementar de avaliação é a \textbf{célula}: o cruzamento de um braço textual, uma
parametrização, um algoritmo de agrupamento e um alfabeto. Dentro de uma célula, cada família de
modelo é ajustada \emph{separadamente em cada grupo} e pontuada sobre os traços retidos daquele
mesmo grupo. Um processo administrativo entra numa célula se satisfaz três condições encadeadas:
possuir traço com ao menos dois estados no alfabeto daquela célula; possuir rótulo naquela partição;
e pertencer a um grupo com ao menos trinta traços. O piso de trinta existe porque um grupo dividido
em 70/30 precisa de casos de avaliação suficientes para que uma acurácia signifique alguma coisa.

O número de grupos avaliados por uma célula difere do $k$ do agrupamento que a originou. Grupos
abaixo do piso são excluídos da avaliação sem serem excluídos da partição, e a primeira condição depende do
alfabeto --- de modo que a mesma partição avalia populações distintas em alfabetos distintos.

\subsubsection*{A divisão treino/avaliação}

Os processos administrativos são divididos em 70/30 entre conjunto de treino e conjunto de avaliação. A
divisão é função determinística do identificador do processo, e não da sua posição em uma lista, de
modo que um processo cai do mesmo lado da divisão independentemente do subconjunto de que seja
extraído --- propriedade que o controle de modelo global, adiante, exige, pois de outro modo um
processo retido do seu próprio grupo poderia comparecer ao conjunto de treino agregado.

\subsubsection*{As medidas, e o eixo que cada uma exige}

A medida primária é a \textbf{acurácia de próximo passo}, calculada sobre \emph{passos de decisão}
e não sobre traços: dado o prefixo de um traço retido, qual é o estado seguinte. Um traço de $n$
estados contribui com $n-1$ passos, de modo que traços longos pesam mais dentro do seu grupo.
Reporta-se a primeira escolha e as três melhores. Ao lado dela mede-se a qualidade da \emph{geração}
por uma distância de edição com transposições entre o sufixo gerado e o sufixo real, condicionada ao
prefixo verdadeiro de metade do traço --- o modelo recebe o que de fato ocorreu até ali e completa o
restante, arranjo em que valores menores são melhores.

Acurácia elevada de próxima atividade não demonstra, por si só, que a estrutura do processo tenha
sido capturada \cite{peeperkorn2023rnnstructure}. Por isso são reportadas, ao lado da acurácia,
medidas de sustentação de caminho: a fração das unidades que um processo real visita e que o
caminho gerado também visita; a fração das unidades do caminho gerado que processos reais visitam;
o número médio de ocorrências de cada transição gerada por membro do grupo; e a contagem de
transições que não ocorrem em \emph{nenhum} membro amostrado do grupo --- fabricação, e não
variação. Mede-se ainda a \textbf{validade estrutural} das gerações completas, isto é, o grau de
conformidade às regras de movimentação que a decodificação da seção~\ref{sec:caminhos} impõe. Três
medidas descrevem a partição e não o modelo, e por isso assumem um único valor por célula, o mesmo
para as seis famílias: a fração das transições do conjunto de avaliação já presentes no grafo de
precedência do conjunto de treino, a sua conversa, e a média harmônica das duas.

Duas ressalvas de alcance acompanham essas definições. A sustentação de transição conta
ocorrências, e não membros distintos, de modo que pode exceder 1 quando uma transição se repete
dentro de um mesmo processo; não é, portanto, uma proporção. As transições não observadas são
contadas contra uma amostra de no máximo trinta traços retidos por grupo, e não contra o grupo
inteiro. O Apêndice A (Tabela~\ref{tab:glossary}) define todas as medidas utilizadas.

Uma regra governa quais dessas medidas cada célula emite, e ela decorre diretamente dos alfabetos
da seção~\ref{sec:alfabetos}: \textbf{uma medida é calculada quando, e somente quando, o eixo que
ela mede existe no alfabeto daquela célula}. O alfabeto de ação por unidade carrega os dois eixos e
emite tudo. O de procedimento carrega apenas a ação, logo emite validade estrutural e não emite
medida alguma de unidade. O de roteamento carrega apenas a unidade, logo emite as medidas de
unidade e \emph{não} emite validade estrutural. Onde o eixo não existe, a chave simplesmente não
aparece no artefato. A alternativa --- calcular a medida assim mesmo --- não devolve um valor
errado, e sim um valor \emph{vazio de referente} que se renderiza como legítimo: uma restrição de
movimentação aplicada a estados que não carregam ação nunca prende, e reportaria conformidade
perfeita por construção.

São três as linhas de base, cada uma respondendo a uma pergunta diferente. A linha de base do
\textbf{sucessor intragrupo} prediz o próximo estado mais frequente dentro do grupo, e pergunta se
o modelo aprendeu alguma coisa. O \textbf{modelo global} é uma única cadeia de Markov de terceira
ordem treinada sobre o conjunto de treino agregado do corpus inteiro e avaliada sobre os mesmos
traços retidos; ele pergunta se um único modelo teria ido igualmente bem. A \textbf{partição
embaralhada} permuta a filiação aos grupos preservando a quantidade de grupos, a distribuição de
tamanhos e a população avaliada; ela pergunta se é o agrupamento, e não a forma da partição, que
está fazendo o trabalho. A esses três controles quantitativos soma-se, na seção 4.4.3, um confronto
qualitativo do caminho gerado contra um processo administrativo real e integralmente conhecido; ele
não substitui nenhuma das linhas de base acima, e serve para expor o que as medidas agregadas não
tornam visível.

Avaliar um agrupamento pelo modelo que ele habilita jusante é princípio consolidado
\cite{deweerdt2013actitrac}. Aqui ele é estendido da aderência do modelo de processo para a predição
de próxima atividade, aplicado externamente a partições produzidas de forma independente --- e não
dentro de um invólucro de agrupamento --- e acrescido dos dois últimos controles, que essa
literatura não reporta rotineiramente. A comparação contra o modelo global é feita com a família de modelos fixada dos dois lados --- cadeia de Markov por grupo contra cadeia de Markov global ---, e não tomando o melhor de seis famílias em cada grupo contra uma única família global. As duas formas são calculadas e reportadas em separado na seção~\ref{sec:res-sinal}, para que o leitor compare o que cada escolha de comparador produz.

\subsubsection*{Agregação: duas leituras, e o que separa uma da outra}

Cada grupo produz a sua própria acurácia; a célula agrega essas acurácias de duas maneiras, e ambas
são reportadas. A leitura \textbf{macro} é a média simples sobre os grupos, e nela cada \emph{tipo de
processo} vale o mesmo, por maior ou menor que seja. A leitura \textbf{ponderada} pesa cada grupo
pelo seu número de traços, e nela cada \emph{processo administrativo} vale o mesmo. Não são duas
versões da mesma grandeza: respondem a perguntas administrativas diferentes --- quão bem o método
serve o catálogo de tipos, e quão bem serve o acervo de casos.

A leitura ponderada carrega uma assimetria: o peso é o número de traços do grupo,
enquanto a acurácia de cada grupo é calculada sobre \emph{passos de decisão}. Um grupo de traços
longos contribui com muitos passos para a sua própria acurácia, mas entra na agregação apenas com o
seu número de processos. Onde a distribuição de comprimentos variar muito entre grupos, as duas
leituras divergem, e a divergência entre elas é o sinal de que isso está ocorrendo.

\subsection{O estimando: ganho contra a partição-controle}\label{sec:estimando}

A grandeza estimada é definida antes de qualquer medição.

\subsubsection*{O dividendo da partição}

Ajustar um modelo por grupo, em vez de um modelo para todo o acervo, altera aquilo que cada modelo
enfrenta. O alfabeto efetivo de um grupo é um subconjunto do alfabeto global, e prever o próximo
passo dentro de um alfabeto menor é mais fácil; em contrapartida, cada modelo passa a dispor de uma
fração dos dados de treino. Nenhum desses dois efeitos depende de a partição significar coisa
alguma: ambos decorrem apenas de quantos grupos existem e de que tamanhos são. Chame-se
\textbf{dividendo da partição} ao efeito líquido dessa mudança de forma.

Um ganho medido contra um modelo global mistura, portanto, duas coisas: o dividendo da partição e o
que a pertença aos grupos de fato carrega. Como o sinal do dividendo não é conhecido \emph{a
priori} --- ele depende de qual dos dois efeitos predomina no acervo em questão --- não é possível
sequer afirmar que um ganho bruto seja um limite inferior do efeito de interesse.

\subsubsection*{A decomposição e o estimando}

Seja $A(\Pi)$ a acurácia agregada dos modelos por grupo sob a partição real $\Pi$, e
$A(\tilde{\Pi})$ a mesma grandeza sob a partição-controle embaralhada da seção~\ref{sec:controles},
que preserva quantidade de grupos, distribuição de tamanhos e população. O ganho bruto contra o
modelo global $B$ decompõe-se em

\[
\underbrace{A(\Pi) - B}_{\text{ganho bruto}} \;=\;
\underbrace{\left[A(\tilde{\Pi}) - B\right]}_{\text{dividendo da partição}} \;+\;
\underbrace{\left[A(\Pi) - A(\tilde{\Pi})\right]}_{\text{sinal da partição}} .
\]

O \textbf{estimando deste trabalho é o segundo termo}. Para uma família de modelo $m$ e uma partição
$\Pi$, com $A_w$ denotando a acurácia agregada ponderada pelo número de traços de cada grupo,

\[
\boxed{\;\Delta^{*}(m,\Pi) \;=\; A_w\!\left(m,\Pi\right) \;-\; A_w\!\left(m,\tilde{\Pi}\right)\;}
\]

\noindent e a variante ajustada, que subtrai o modelo global sob ambas as partições,

\[
\Delta^{*}_{\mathrm{aj}}(m,\Pi) \;=\;
\left[A_w(m,\Pi) - B(\Pi)\right] \;-\; \left[A_w(m,\tilde{\Pi}) - B(\tilde{\Pi})\right],
\qquad
r \;=\; B(\Pi) - B(\tilde{\Pi}) \;\approx\; 0 ,
\]

\noindent em que $r$ é o resíduo: o mesmo modelo global pontuado sobre os mesmos passos retidos sob
as duas partições, que difere de zero apenas pelo efeito da ponderação e por isso serve de teste de
sanidade da agregação. $\Delta^{*}$ é, portanto, a diferença entre a partição real e um
embaralhamento de forma idêntica. É o único dos dois do qual o dividendo foi
subtraído, e é positivo apenas se a pertença aos grupos carregar informação que a forma da partição
não carrega.

\subsubsection*{$\Delta^{*}$ é portão de admissão, não ordenação}

O controle preserva a forma da partição, mas não equaliza o \emph{custo
marginal} de fragmentar os dados de treino: um grupo formado ao acaso tem dados de treino sem
estrutura interna, e o seu modelo degrada-se com a redução de tamanho de modo distinto do de um
grupo coerente. Logo, $\Delta^{*}$ não é invariante à granularidade e não pode ser usado para
ordenar partições com números de grupos diferentes.

O uso legítimo é como \textbf{critério de admissão}: uma partição cujo $\Delta^{*}$ não seja
positivo não carrega sinal além da sua forma e não deve integrar o estudo, qualquer que seja a sua
validade interna. Para \emph{escolher} entre partições admitidas emprega-se o critério de seleção da
seção~\ref{sec:desenho}, que pondera desempenho a jusante, cobertura e qualidade. Comparações a
número de grupos emparelhado permanecem válidas sob $\Delta^{*}$.

\subsubsection*{Os dois comparadores respondem a perguntas diferentes}

Nem o modelo global nem o embaralhamento é dispensável, porque as perguntas não são a mesma. O
modelo global pergunta \emph{um único modelo teria ido igualmente bem?} --- uma questão de
engenharia, cuja resposta orienta quem decide se vale a pena manter um modelo por tipo. O
embaralhamento pergunta \emph{este agrupamento significa alguma coisa?} --- uma questão de validade,
cuja resposta orienta quem decide se o catálogo é um objeto real. Os dois são calculados e
reportados lado a lado, por alfabeto, e nenhum é apresentado como substituto do outro.

\subsection{A camada de evidência: o que cada verificação tenta quebrar}\label{sec:evidencia}

As etapas descritas até aqui produzem o artefato. As etapas seguintes o testam, e cada uma responde
a uma objeção específica. Todas são executadas por alfabeto, quando a grandeza que medem depender do alfabeto.

\begin{itemize}
  \item \textbf{Remoção controlada de blocos e de histórico} --- \emph{é o conteúdo documental que agrupa, ou
    qualquer coisa agruparia?} Blocos de atributos são zerados um a um e o histórico de cada
    processo é truncado, medindo-se quanto a partição se desloca.
  \item \textbf{Viés do corpus} --- \emph{os processos que concluem dentro do ano representam os
    que foram abertos nele?} A censura à direita é estimada mês a mês (seção~\ref{sec:vies}).
  \item \textbf{Retenção temporal} --- \emph{a vantagem sobrevive quando a partição vê apenas parte
    da história de cada caso?} Os cortes são reavaliados sobre os casos que todos eles avaliam,
    porque de outro modo comparar-se-iam populações distintas.
  \item \textbf{Conformidade} --- \emph{quanto do comportamento observado o modelo descoberto
    reproduz, e quanto a mais ele permite?}
  \item \textbf{Custo da restrição} --- \emph{o que a garantia de validade cobra do gerador?}
  \item \textbf{Não inferioridade} --- \emph{a supressão de dados pessoais custa utilidade, contra
    uma margem fixada antes de medir?} Reamostragem pareada por processo, por alfabeto.
  \item \textbf{Estabilidade} --- \emph{a partição sobrevive à reamostragem dos seus próprios
    membros?}
  \item \textbf{Auditoria de reajuste} --- \emph{cada agrupamento armazenado se reproduz a partir da
    matriz de atributos que ele próprio declara?} A verificação reexecuta a receita declarada, o que
    inclui a política de ruído da seção~\ref{sec:ruido}.
  \item \textbf{Partição-controle} --- introduzida na seção~\ref{sec:controles} e formalizada na
    seção~\ref{sec:estimando}, é a que decide se qualquer das anteriores tem sobre o que incidir.
\end{itemize}

\textbf{O comparador decide a resposta}: os mesmos modelos, sobre os mesmos casos, produzem
afirmações distintas conforme aquilo contra o que são medidos. O estimando é fixado antes de medir,
e o comparador é reportado ao lado de cada número.

\subsection{Desenho experimental, reprodutibilidade e verificação}\label{sec:desenho}

\textbf{Braços e células.} O \emph{pipeline} é executado duas vezes, produzindo dois braços que
diferem em uma única coisa: o texto de entrada. O braço \textbf{sem anonimização} usa o documento
inicial como ele consta do sistema; o braço \textbf{anonimizado} usa o mesmo documento após a etapa de
pseudonimização e anonimização da seção~\ref{sec:documento}. Todo o restante --- blocos de atributos,
parâmetros, sementes, protocolo de avaliação --- é idêntico, o que torna a diferença entre os
braços atribuível à anonimização e não a outra escolha.

Cada braço produz \textbf{sete} agrupamentos, resultado de três famílias algorítmicas em
configuração de referência mais quatro variantes de parâmetro:

\begin{itemize}
  \item \textbf{HDBSCAN} --- agrupamento por densidade, que não exige número de grupos declarado; a
    variante \emph{hcoarse} funde grupos separados por menos que um limiar de distância, produzindo
    partição mais grosseira.
  \item \textbf{k-means} --- agrupamento por centroide, com $k$ declarado; entram duas variantes,
    uma com $k$ escolhido pelo método do cotovelo (\emph{kelbow}) e outra por \emph{silhouette}
    (\emph{ksil}).
  \item \textbf{Leiden} --- agrupamento por modularidade sobre grafo de vizinhança; a variante
    \emph{lres} altera o parâmetro de resolução.
\end{itemize}

\noindent São, portanto, \textbf{catorze} agrupamentos reais --- sete por braço ---, cada um com a
sua própria partição-controle embaralhada (seção~\ref{sec:controles}), o que perfaz \textbf{vinte e
oito} partições. Cada uma é avaliada nos \textbf{três} alfabetos de traço da
seção~\ref{sec:alfabetos}, totalizando \textbf{84} células de avaliação.

Dentro de cada célula, as \textbf{seis} famílias de gerador de caminho típico da
seção~\ref{sec:caminhos} são treinadas por grupo, acima do piso de trinta traços.

\textbf{Controles.} O controle de partição embaralhada é gerado por permutação dos rótulos de um
agrupamento real, o que preserva a quantidade de grupos, a distribuição de tamanhos e a população
que sobrevive ao piso de avaliação; a única diferença em relação à partição real é, portanto, a
filiação. A lista das células de controle é derivada da lista das células reais pela própria etapa
que as executa, e não escrita em paralelo a ela, de modo que as duas não possam divergir sem
falhar.

\textbf{A escolha entre partições admitidas.} $\Delta^{*}$ decide quais partições entram no estudo
(seção~\ref{sec:estimando}), mas não as ordena. O critério que as ordena pontua os catorze
agrupamentos sobre quatro grandezas escalonadas por mínimo--máximo entre os candidatos: a acurácia
sobre o \emph{conjunto comum} --- os processos administrativos que \emph{todas} as partições
avaliam, sem o qual se compararia desempenho sobre populações distintas ---, a \textbf{cobertura},
a \emph{silhouette} e o \textbf{tamanho mediano} de grupo. Todo agrupamento é candidato, sem filtro
prévio de cobertura, uma vez que atribuir o ruído da densidade ao grupo mais próximo
(seção~\ref{sec:ruido}) já fechou o canal de descarte a montante. O escore é calculado sob quatro
ponderações, uma delas ignorando por completo os índices de validade, e apenas no alfabeto de
roteamento --- os alfabetos nunca são agregados entre si, de modo que aquele sobre o qual se
seleciona é decisão de escopo, e não propriedade dos dados. Um candidato que vença sob as quatro
vence sem depender de qual delas se adote.

Um segundo escore, restrito aos seis agrupamentos principais, serve de verificação de sensibilidade:
três termos com pesos declarados de antemão --- \textbf{desempenho a jusante}, com peso 0,5, medido
sobre o conjunto comum e escalado a partir do piso que a melhor partição embaralhada atinge;
\textbf{cobertura}, com peso 0,3, a fração do corpus que sobrevive ao piso de trinta traços; e
\textbf{qualidade interna}, com peso 0,2, a média de dois índices de validade reescalados ---,
varridos sobre cinco esquemas de ponderação: equilibrado, jusante em primeiro, cobertura em
primeiro, qualidade em primeiro e terços iguais. Os dois critérios são reportados lado a lado, com o
vencedor sob cada esquema. Uma escolha que só sobrevive a um esquema de pesos é uma escolha do
esquema.

\textbf{Reprodutibilidade.} Cada etapa declara a pergunta que responde, se é ou não determinística
e, quando não é, por quê; e copia essa declaração para dentro de cada artefato que escreve, ao lado
do identificador da versão do código, das versões das bibliotecas numéricas empregadas, dos
parâmetros resolvidos e do resumo criptográfico de cada entrada consumida. As matrizes de atributos
carregam o mesmo carimbo, e as etapas de agrupamento e de avaliação recusam-se a executar contra
atributos cujo carimbo não corresponda ao braço que declaram --- guarda, e não convenção. Uma etapa
de auditoria reajusta cada agrupamento armazenado a partir da sua própria matriz carimbada, replica
a receita que ele declara, inclusive a política de ruído da seção~\ref{sec:ruido}, e compara o
resultado aos rótulos armazenados; o desfecho dessa auditoria é reportado na
seção~\ref{sec:res-entrada}. Uma segunda auditoria percorre a árvore inteira de artefatos e verifica
que cada arquivo ainda declara a etapa que o produziu e que as entradas declaradas continuam a
resumir do mesmo modo. É esse encadeamento --- pergunta, versão, entradas, reajuste --- que
converte ``isto foi executado'' em ``qualquer pessoa pode reexecutar isto''.

As vetorizações de documento e de tópico são mantidas em cache em disco, e assim alternar entre
braços restaura vetores idênticos em vez de recomputá-los. Um documento cujo texto anonimizado seja
idêntico ao bruto conserva o vetor que já possui, em vez de ser reprocessado; ele não é excluído,
mas tampouco contribui com sinal de anonimização, e a etapa registra quantos documentos estão nessa
condição.

\textbf{Dos artefatos a este texto.} Toda tabela quantitativa deste trabalho é gerada a partir de
artefatos armazenados por uma única etapa do pipeline, que também regenera as tabelas
correspondentes no documento interno de resultados do projeto. Um modo de verificação re-renderiza
tanto as tabelas do manuscrito quanto as do documento interno e falha se qualquer uma delas houver
divergido dos artefatos, de modo que uma cifra não pode ser editada dentro de uma tabela sem que a
evidência subjacente mude primeiro. A grade completa de combinações de parâmetros não é impressa
aqui; ela é regerável a partir dos artefatos armazenados.

Uma verificação de regeneração alcança as tabelas geradas, mas não o texto corrido, que é
escrito à mão. Mantém-se, por isso, uma segunda verificação, executada com a suíte de testes:
um registro explícito das grandezas que sustentam o argumento na prosa, que resolve cada uma
até o seu caminho na árvore de artefatos e falha nos dois sentidos --- se uma cifra do texto
deixa de arredondar para o artefato que cita, ou se o artefato muda e a frase não é atualizada
com ele. O registro cobre as duas versões do manuscrito e inclui asserções de sinal, de modo
que uma redação que preserve a magnitude mas apague o sinal do resultado seja reprovada. Ele não é
exaustivo: as cifras que ele cobre não podem derivar silenciosamente; qualquer cifra fora dele
está verificada apenas pela leitura.

\subsection{Procedimentos estatísticos}\label{sec:estatistica}

A comparação entre os braços sem anonimização e anonimizado emprega \emph{bootstrap} pareado sobre os processos
administrativos que os dois braços avaliam, com 2.000 replicações, e é conduzida por alfabeto,
jamais entre alfabetos. O par é o processo, e não a célula, o que remove do intervalo a variação
que se deve à dificuldade do caso. A conclusão de não inferioridade exige que o limite inferior
unilateral de 90\% da diferença supere uma margem de $0{,}01$ ponto de acurácia, registrada antes
de os dados existirem, e uma margem de sensibilidade de $0{,}005$ é reportada ao lado, de modo que
a robustez da conclusão à escolha da margem fique visível em vez de precisar ser suposta.

A relação entre a qualidade interna de uma partição e a sua qualidade a jusante é testada por
correlação de postos, sob teste de permutação de Monte Carlo com 20.000 sorteios, sobre as catorze
parametrizações reais de um mesmo alfabeto. São testadas oito grandezas: quatro índices internos
de validade --- DBCV, \emph{silhouette}, Davies--Bouldin e Calinski--Harabasz --- e quatro
descritores da forma da partição --- número de grupos, tamanho mediano de grupo, taxa de ruído
final e fração que a densidade recusaria antes da atribuição da seção~\ref{sec:ruido}. Todas são
testadas sem correção para multiplicidade, e todas exigem que os índices provenham de um \emph{mesmo
espaço métrico}: índices calculados em espaços distintos não são comparáveis e não podem ser
reunidos numa única correlação, pelo que o procedimento interrompe a execução em vez de agregar
silenciosamente uma partição fora do espaço comum.

A regra de enunciado desse teste é fixada antes de o resultado ser conhecido: na ausência de
rejeição, reporta-se ausência de evidência, e não se formula afirmação de equivalência a respeito
dos índices. O procedimento, sozinho, não estabelece nem uma coisa nem a outra. O poder do teste,
dado o número de observações disponíveis, é quantificado ao lado do resultado na
seção~\ref{sec:res-indices}.

\subsection{Aspectos éticos, instrumento autorizador e disponibilidade dos
dados}\label{sec:etica}

\textbf{Financiamento.} Nenhuma bolsa ou financiamento externo apoiou este trabalho. A
administração participante prestou apoio institucional em espécie: acesso ao seu \emph{data lake} e
a servidores que explicaram o significado das unidades organizacionais e dos tipos de ação.

\textbf{Ética.} Não foram coletados dados de participantes humanos, pelas razões detalhadas na
seção~\ref{sec:caracterizacao}. Sobre essa base, o estudo não foi submetido a comitê de ética em
pesquisa. A consulta informal a servidores sobre o significado de unidades organizacionais e de
tipos de ação é divulgada aqui.

\textbf{Instrumento autorizador.} O estudo foi conduzido sob autorização formal de tratamento de
dados, expedida pelo titular da administração participante na condição de controladora, nos termos
da Lei Geral de Proteção de Dados \cite{brasil2018lgpd}, e vigente para o período que abrange este
trabalho \cite{seadpi2026termo}. O instrumento autoriza o tratamento de registros de processos
administrativos para fins acadêmicos (cláusula 2.1); exige pseudonimização ou anonimização antes
das fases de clusterização e vetorização (cláusula 5.1); exige que o tratamento ocorra inteiramente
dentro do ambiente computacional da administração, com modelos de pesos abertos executados
localmente e sem transferência a serviços externos (cláusulas 4.1 e 4.5); e restringe a divulgação
acadêmica a nomes de grupos, exemplos genéricos ou sintéticos, métricas agregadas e código-base sem
segredos (cláusula 4.4.1). Todos os resultados aqui reportados situam-se dentro desse escopo. O
Anexo A reproduz o documento na íntegra, com tarja sobre os identificadores pessoais.

Duas consequências desse instrumento incidem diretamente sobre a forma deste trabalho e devem ser
lidas em conjunto com a declaração de posicionalidade da seção~\ref{sec:caracterizacao}. A primeira é a cláusula 4.4.1,
que submete a divulgação acadêmica à aprovação prévia da SEAD-PI, na qualidade de camada adicional
de controle; essa aprovação foi solicitada e concedida para este trabalho. A segunda é a cláusula
7.2, que atribui à Secretaria a titularidade da base de dados e do grafo derivado, com licença de
uso acadêmico ao servidor --- o que restringe, e não meramente organiza, o que pode ser
disponibilizado.

\textbf{Base legal do tratamento.} A autorização descrita acima é o instrumento sob o qual este
estudo transcorreu, mas uma autorização expedida por controlador não é, ela própria, base legal de
tratamento sob a lei geral de proteção de dados. O instrumento aloca responsabilidade pela etapa
que importa: dispõe que o grafo inicial de exploração poderá conter dados não anonimizados, que
essa etapa integra o tratamento sob responsabilidade da controladora, restrita ao ambiente seguro,
e que não constitui uso acadêmico autônomo (cláusula 5.2, transcrita no Anexo A). O que ele não faz
é identificar qual hipótese legal a controladora invocou --- e este trabalho tampouco o faz. Uma administração que pretenda reproduzir este trabalho
precisaria estabelecer essa base antes de começar, e a seção 4.10 expõe por que essa é a primeira
pergunta que tal administração deveria responder.

\textbf{Disponibilidade dos dados.} Sob o instrumento vigente, a cláusula 7.1 atribui a mim a
metodologia e o código-base, ao passo que a base de dados e o grafo derivado são da
Secretaria (cláusula 7.2). Podem, por isso, ser disponibilizados o código-fonte e os artefatos
agregados derivados --- métricas em nível de grupo, as tabelas de resultados geradas e o pipeline
de avaliação. Não podem sê-lo os registros de atividades subjacentes, os documentos dos processos
administrativos e a base de grafo: contêm dados pessoais, e o instrumento veda a sua transferência
a terceiros. Toda afirmação quantitativa deste trabalho é
regerada de artefatos armazenados por uma única etapa do pipeline, que falha se o manuscrito e os
artefatos divergirem.

\textbf{Uso de ferramentas de inteligência artificial.} Modelos de linguagem de grande porte foram
utilizados em dois papéis distintos, ambos aqui divulgados; cada um suscita questões diferentes.

Como \emph{componentes do método}: um modelo de visão computacional executado localmente realizou a
leitura estruturada dos documentos iniciais; um extrator neural de entidades nomeadas
\cite{zaratiana2025gliner2} participou da detecção de dados pessoais descrita na seção~\ref{sec:documento}; um
modelo de linguagem de pesos abertos produziu o tópico e o resumo de cada documento inicial, o
vocabulário controlado de palavras-chave e os nomes dos tipos de processo da seção~\ref{sec:nomeacao}; e um modelo
de vetorização produziu os vetores de documento. Todos foram executados na infraestrutura do
Estado, como exposto na seção~\ref{sec:atributos} --- exigência do instrumento autorizador, não preferência técnica
---, e nenhum dado de processo administrativo foi enviado a qualquer serviço externo. Esses
componentes são objeto de avaliação, e não de confiança: a seção 4.11 registra que os nomes gerados
são saída de máquina não validada.

Como \emph{assistência de revisão adversarial}: houve emprego de assistência de inteligência
artificial na revisão do texto do manuscrito e do código de análise e de testes, em especial na
identificação de lacunas lógicas --- afirmações não sustentadas pela evidência citada, números
divergentes entre prosa e artefato, e passos do argumento omitidos. Esse uso é distinto do
anterior: não produziu texto, e sim objeções a ele. Diversas correções registradas neste trabalho
--- entre elas a retratação de resultados descrita na seção 4.11 --- originaram-se dessa revisão.
Avaliei cada objeção individualmente, e as acolhidas foram conferidas contra os artefatos.

Conferi cada cifra reportada contra os artefatos gerados e respondo por todas as afirmações aqui
feitas.
